\scanpage{70}
\deleted{\uline{Exemples} 7.5.3. \begin{tabular}{@{}c@{}}\inserted{Chacune des}\\\deleted{Les}\end{tabular} propriétés \raisebox{1.7ex}{\inserted{\uuline{R}}} suivantes satisfait aux \begin{tabular}{@{}l@{}}\hspace{1.6em}\figure{\begin{tikzpicture}[baseline=0pt]\draw[line width=.4pt] (0,0) -- (-.01,.13);\end{tikzpicture}}\\condition\end{tabular} a) et b) de 7.5.1., \deleted{pourvu que $B$ soit un} lorsque $B$ est un excellent anneau:
\begin{enumerate}
\item[(i)] $\operatorname{coprof}C\quad n$
\item[(ii)] $S_n$
\item[(iii)] Coh Mac
\item[(iv)] régulier
\item[(v)] normal $R_n$
\item[(vi)] normal
\item[(vii)] intègre
\item[(viii)] réduit
\end{enumerate}

En fait, pour (vii) et (viii) c'est vrai sans aucune hypothèse sur $B$, pour (i) (ii) (iii) il suffit de \begin{tabular}{@{}l@{}}\inserted{et la cond univ cat dans (ii)}\\vérifier la condition (b), \deleted{qui}\end{tabular} ce \begin{tabular}{@{}l@{}}\inserted{qui}\\est\end{tabular} vérifiéé si \deleted{\unsure{tout|tant} que} $B$ est quotient d'un régulier, ou satisfait la condition de Nagata sur le \begin{tabular}{@{}l@{}}\inserted{(et dans le cas (ii) est univ cat).}\\lieu singulier. Cette condition est\end{tabular} également suffisante pour faire marcher (iv) (v) (vi), pourvu que dans (v) et (vi) on suppose de plus $B$ universellement caté\deleted{\unsure{h|l}}naire. Enfin, pour faire marcher (vi), on peut aussi prendre la condition ``$B$ est univ japonais et univ cat''.
\margin{gata\figure{\begin{tikzpicture}[baseline=0pt]\draw[line width=.35pt,opacity=.65] (-.05,.10) .. controls (.09,.11) and (.10,.03) .. (.02,-.01);\fill[opacity=.65] (.04,-.05) circle (.015cm);\end{tikzpicture}}}

\uline{Remarques} 7.5.4.\quad (i) comme 7.5.6. (i)

(ii) Comme 7.5.6. (ii) ; il s'agirait pratiquement des propriétés (i) \deleted{(ii)} (iii) (vii) ou (viii).

(iii) Dans la situation 7.5.1., on dispose le plus souvent d'une propriété supplémentaire de \uuline{R}, garantissant que si \deleted{$A\quad B$} $\text{\uuline{R}}(A)$ est vrai et les fibres de $\operatorname{Spec}(B)\quad\operatorname{Spec}(A)$ satisfont la condition \uuline{R} géométrique, \deleted{il en est de même} alors \deleted{($\operatorname{Spec}(B)$ étant} ($B$ étant plat sur $A$, rappelons-le) $\text{\uuline{R}}(B)$ est vrai. On peut donc ti\deleted{e}rer de 7.5.1. que}
